\title{Controlling Text-to-Image Diffusion by Orthogonal Finetuning}

\begin{abstract}
Text-to-image diffusion models at large scales demonstrate remarkable proficiency in creating photorealistic visuals from textual descriptions. However, efficiently steering these sophisticated models to solve a variety of downstream tasks remains an unresolved question. To address this, we propose a systematic finetuning approach termed Orthogonal Finetuning~(OFT) for adapting text-to-image diffusion models. Distinct from previous techniques, OFT offers a theoretical guarantee of maintaining the hyperspherical energy, a metric quantifying the relationships between neurons constrained on a unit hypersphere. We observe that retaining this property is essential for upholding the semantic generative capacity of text-to-image diffusion models. To enhance stability during finetuning, we introduce Constrained Orthogonal Finetuning (COFT), which incorporates an explicit constraint on the hypersphere’s radius. Our focus is on two primary finetuning applications: subject-driven generation, where the model synthesizes subject-specific images based on limited examples and a textual cue, and controllable generation, where the model is conditioned by supplementary signals. Through empirical evaluations, we demonstrate that OFT delivers superior results over existing alternatives, with gains in both generative quality and training efficiency.
\end{abstract}

\section{Introduction}

State-of-the-art text-to-image diffusion models~\cite{saharia2022photorealistic,ramesh2022hierarchical,rombach2022high} have set new standards in text-guided image synthesis, excelling at producing high-quality images. Nonetheless, textual instructions may lack the precision and granularity needed for detailed or accurate image control. To mitigate this limitation, we explore two categories of text-to-image generation problems in this work:

\begin{itemize}[leftmargin=*,nosep]

    \item \textbf{Subject-driven generation}~\cite{ruiz2023dreambooth}: Provided with a handful of images depicting a particular subject, this task involves creating new images of the same subject in varied scenarios based on a textual description.
    \item \textbf{Controllable generation}~\cite{zhang2023adding,mou2023t2i}: Leveraging an extra control input (\emph{e.g.}, canny edges, segmentation maps), the goal here is to synthesize images that are aligned with both the specified control and the text prompt.
\end{itemize}

The central challenge in both tasks is to adapt text-to-image diffusion models for specific purposes while maintaining the strong generative capabilities acquired during pretraining. We identify two primary criteria for an effective finetuning approach: (1) \emph{training efficiency}, referring to minimizing the number of tunable parameters and required training epochs, and (2) \emph{generalizability preservation}, which entails retaining the pretrained model's high-quality and diverse generation. Conventional finetuning methodologies typically involve either making minor adjustments to existing neuron weights with a low learning rate (\emph{e.g.}, \cite{ruiz2023dreambooth}), or incorporating small, re-parameterized components into the network (\emph{e.g.}, \cite{hulora2022,zhang2023adding}). However, despite their straightforwardness, these techniques do not ensure that the model’s generative performance, as established in pretraining, will be preserved. Furthermore, there is a notable absence of a systematic framework for developing effective finetuning schemes or for selecting key hyperparameters, such as the optimal number of training epochs. A fundamental obstacle is the absence of a metric to accurately assess the retention of the pretrained generative capabilities. Prevailing finetuning strategies generally rely on the implicit assumption that reducing the Euclidean distance between the weights of the finetuned and pretrained models leads to better preservation of the original capabilities. Consequently, very small learning rates are often employed. While this assumption may sometimes hold true, we contend that the Euclidean distance alone inadequately reflects the preservation of semantic information. Thus, developing a more structurally-informed metric to evaluate the distinctions between finetuned and pretrained models could significantly enhance both the fidelity of pretraining performance retention and the overall stability of the finetuning process.

Building upon empirical findings that hyperspherical similarity serves as an effective means of representing semantic relationships~\cite{liu2017deep,liu2018decoupled,chen2020angular}, we utilize hyperspherical energy~\cite{liu2018learning} to describe the pairwise relational configuration of neurons. Specifically, hyperspherical energy, computed as the aggregate hyperspherical similarity (\emph{e.g.}, cosine similarity) across all pairs of neurons within the same layer, quantifies how uniformly the neurons are distributed over the unit hypersphere~\cite{liu2021learning}. We conjecture that successful finetuning should result in minimal alteration of the hyperspherical energy relative to the original pretrained model. A straightforward approach would be to regularize the training so as to maintain constant hyperspherical energy throughout finetuning; however, such a method does not ensure that discrepancies in hyperspherical energy will be adequately constrained. To address this, we leverage an invariance property: pairwise hyperspherical similarities are rigorously maintained under the application of a shared orthogonal transformation to all neurons. Drawing on this invariance, we introduce Orthogonal Finetuning~(OFT), a technique that adapts large-scale text-to-image diffusion models for downstream tasks while ensuring that hyperspherical energy remains unchanged. The core concept is to learn an orthogonal transformation, common across layers, that maintains the relative angles among neurons. OFT can alternatively be interpreted as redefining the reference coordinate system for neurons within each layer. Furthermore, taking into account that a reduced Euclidean distance between the finetuned and pretrained models suggests superior retention of the original model's performance, we put forward an extension: Constrained Orthogonal Finetuning~(COFT), which restricts the adapted model to remain within a hypersphere of fixed radius centered at the pretrained neuron positions.

The motivation for utilizing orthogonal transformations in the finetuning of neurons partly stems from concepts in the 2D Fourier transform. Specifically, a 2D Fourier transform enables an image to be represented in terms of magnitude and phase spectra. Importantly, the phase spectrum—which encapsulates the angular relation between the input and the basis functions—retains most of the image's semantic content. As an illustration, even if the magnitude spectrum of an image is replaced with a random one, the image reconstructed using its original phase spectrum preserves its essential semantics. This observation implies that adjusting neuron directions is central to introducing semantic changes to generated images, which is a primary objective in subject-driven and controllable generation settings. However, allowing neuron directions to change without constraint can severely degrade the generative abilities conferred by pretraining. To address this, we advocate for maintaining the relative angles between all pairs of neurons, guided by the premise that such angles encode critical knowledge within neural networks. This perspective leads us to propose a layer-wise orthogonal transformation—shared across each layer—that ensures the hyperspherical energy remains invariant.

Our work is also informed by techniques from orthogonal over-parameterized training~\cite{liu2021orthogonal}, which optimize classification networks by applying orthogonal transformations to networks initialized at random. Since random initialization yields a theoretically low expected hyperspherical energy, \cite{liu2021orthogonal} aims to sustain this property throughout training, with evidence indicating that lower energy enhances generalization in classification tasks~\cite{liu2018learning,lin2020regularizing}. The findings of \cite{liu2021orthogonal} suggest that orthogonal transformations offer sufficient capacity for learning neural networks that generalize well to classification scenarios. In contrast, our investigation centers on the finetuning of text-to-image diffusion models, targeting both enhanced control and improved generative capacity on downstream tasks. We draw a distinction between OFT and the approach in \cite{liu2021orthogonal} on two primary grounds. First, whereas \cite{liu2021orthogonal} is tailored toward minimizing hyperspherical energy, OFT is constructed to conserve the hyperspherical energy of the original, pretrained model so that its inherent structures are preserved post-finetuning; minimizing this energy during finetuning, particularly in diffusion models, risks eliminating important semantic content. Second, OFT explicitly limits deviations from the pretrained parameters, introducing a constrained approach not present in \cite{liu2021orthogonal}. The essence of effective finetuning lies in achieving an optimal compromise between adaptability and the retention of pretrained knowledge, and we contend that our OFT methodology accomplishes this. Below, we summarize our primary contributions:

\begin{itemize}[leftmargin=*,nosep]

    \item We introduce Orthogonal Finetuning, a new approach for finetuning text-to-image diffusion models that enhances their controllability. To further address stability concerns, we devise a constrained version that restricts the angular divergence from the pretrained model.
    \item In contrast to other finetuning strategies, OFT modifies the model while rigorously maintaining hyperspherical energy, an aspect we empirically establish as critical for preserving the generative semantics of the original pretrained model.
    \item Our evaluation encompasses OFT applied to both subject-driven and controllable generation tasks. Through extensive experiments, we report substantial gains over established techniques with respect to image quality, convergence rate, and the stability of finetuning. Additionally, OFT demonstrates superior sample efficiency, requiring significantly fewer images and epochs to achieve satisfactory results.
    \item For tasks involving controllable generation, we propose a novel control consistency metric that quantifies model controllability. This metric works by extracting the control signal from a generated image and comparing it against the initial control signal, providing further empirical support for the robust controllability achieved by OFT.
\end{itemize}

\section{Related Work}

\textbf{Text-to-image diffusion models}. The field of text-to-image synthesis has seen remarkable advancements~\cite{saharia2022photorealistic,ramesh2022hierarchical,rombach2022high,gu2022vector,nichol2022glide}, primarily due to progress in diffusion-based generative methods~\cite{ho2020denoising,song2021score,song2021denoising,dhariwal2021diffusion} and integrated vision-language models~\cite{su2020vlbert,lu2019vilbert,radford2021learning,wang2021simvlm,li2022blip,li2023blip,alayrac2022flamingo,schuhmann2022laion}. For instance, GLIDE~\cite{nichol2022glide} and Imagen~\cite{saharia2022photorealistic} carry out diffusion in the pixel domain, with GLIDE jointly training a text encoder and diffusion prior on paired datasets, whereas Imagen employs a static, pretrained text encoder. Models such as Stable Diffusion~\cite{rombach2022high} and DALL-E2~\cite{ramesh2022hierarchical} instead operate in a latent space; Stable Diffusion incorporates VQ-GAN~\cite{esser2021taming} to create a latent visual codebook, while DALL-E2 leverages CLIP~\cite{radford2021learning} to realize a joint embedding space for both image and text representations. Beyond diffusion models, other generative paradigms like generative adversarial networks~\cite{reed2016generative,zhang2017stackgan,xu2018attngan,li2019controllable} and autoregressive approaches~\cite{ramesh2021zero,ding2021cogview,wu2022nuwa,yuscaling} have been explored for this task. Importantly, OFT is designed to be compatible with any text-to-image diffusion framework due to its model-agnostic finetuning design.

\textbf{Subject-driven generation}. Methods to maintain the integrity of the subject typically include utilizing user-specified masks as supplemental conditions~\cite{avrahami2022blended,nichol2022glide}. Inversion-based techniques~\cite{choi2021ilvr,dhariwal2021diffusion,ramesh2022hierarchical,gal2022image} permit alteration of context while retaining subject features. Local and global edits can be carried out without explicit masks~\cite{hertz2022prompt}. However, such methods often struggle with preserving fine identity-related details. Pivotal Tuning~\cite{roich2022pivotal} refines a generator around a given inverted latent representation, introducing extra regularization for identity maintenance. Analogously, \cite{nitzan2022mystyle} generates a person-specific face prior from multiple facial images, and \cite{casanova2021instance} creates various instances though occasionally sacrifices instance-specific fidelity. DreamBooth~\cite{ruiz2023dreambooth} enhances the text-to-image diffusion model by introducing a tailored token and a handful of subject-specific images, utilizing both reconstruction and class-aware prior preservation losses during finetuning. Building on the DreamBooth setup, OFT instead employs orthogonal transformations in finetuning rather than straightforward updates with a limited learning rate.

\textbf{Controllable generation}. Image-to-image translation represents a specific instance of controllable generation, with earlier approaches largely relying on conditional generative adversarial networks~\cite{isola2017image,zhu2017toward,wang2018high,park2019semantic,choi2018stargan}. Diffusion models have also emerged as competitive tools for image-to-image translation~\cite{saharia2022palette,voynov2022sketch,wang2022pretraining}. In recent developments, ControlNet~\cite{zhang2023adding} introduces a strategy to exert fine-grained control over a pretrained diffusion model by finetuning it to accommodate auxiliary control inputs, achieving notable advancements in controllable generation. Similarly, the contemporaneous T2I-Adapter~\cite{mou2023t2i} pursues enhanced controllability by finetuning an existing diffusion model. Adopting the experimental design of~\cite{zhang2023adding,mou2023t2i}, our method applies OFT to the finetuning of pretrained diffusion models, consistently producing superior control while requiring less training data and a reduced number of parameters for finetuning. Importantly, OFT does not entail any extra inference-time computational cost.

\textbf{Model finetuning}. The strategy of adapting large pretrained models for specific downstream applications has seen widespread adoption~\cite{devlin2018bert,brown2020language,he2022masked}. Adaptation techniques, such as those in~\cite{hulora2022,houlsby2019parameter,pfeiffer2020adapterfusion}, are especially well-explored in the natural language processing field. OFT is most closely related to LoRA~\cite{hulora2022}, which incorporates a low-rank hypothesis for the parameter update during the finetuning phase. Diverging from this, OFT implements a shared orthogonal transformation across layers to adjust neuron weights multiplicatively. This approach theoretically conserves the angular relationships between neurons within the same layer, contributing to improved stability.

\section{Orthogonal Finetuning}

\subsection{Why Does Orthogonal Transformation Make Sense?}

To motivate the use of orthogonal transformations when finetuning text-to-image diffusion models, we break the reasoning into two central questions: (1) the motivation for adjusting neuron angles (that is, directions) during finetuning, and (2) the rationale for employing orthogonal transformations to effect such angle modifications.

Addressing the first question, our approach is informed by empirical findings reported in \cite{liu2018decoupled,chen2020angular}, which indicate that angular differences among features effectively capture the semantic gap. SphereNet~\cite{liu2017deep} demonstrates that enforcing normalization of all neurons onto a unit hypersphere during neural network training maintains representational capacity and can enhance generalization, suggesting that the orientation of neuron vectors suffices to encode the critical information from the data. To illustrate the centrality of neuron angles, we conduct a controlled experiment as visualized in Figure~\ref{fig:toy_ae}, utilizing a conventional convolutional autoencoder on a dataset of flower images. During training, the typical inner product computes the feature maps (with $z$ representing the output of convolution kernel $\bm{w}$ on input $\bm{x}$ within each sliding window). At test time, we contrast three distinct strategies for feature map computation: (a) the same inner product from training, (b) reliance on magnitude only, and (c) exclusively angular information. As shown in Figure~\ref{fig:toy_ae}, preserving only the angular components of neuron activations allows almost perfect reconstruction of the input images, while the magnitude information alone offers negligible value. It is crucial to note that cosine activation (c) is not employed during training; training proceeds solely using the inner product. These observations suggest that neuron angles (i.e., directions) are the predominant carriers of semantic information within the images, indicating that modifying neuron directions should be prioritized when seeking to alter image semantics.

With regard to the second question, the most straightforward method to adjust neuron directions is to collectively apply a rotation or reflection to all neurons in a layer, which inherently introduces an orthogonal transformation. Although it is conceivable to utilize more general angular transformations, such as assigning distinct rotation angles to each neuron, orthogonal transformations strike an optimal balance between expressive power and structural discipline. Furthermore, prior work \cite{liu2021orthogonal} has established that orthogonal transformations provide ample capacity for neural network learning. To empirically substantiate this point, we execute an experiment in the context of the ControlNet~\cite{zhang2023adding} controllable generation task, with results displayed in Figure~\ref{fig:ortho}. The findings show that standard OFT exhibits stable and precise control by the end of training (epoch 20), whereas eliminating the orthogonality constraint from OFT results in failure to generate plausible images or exercise any controllable effect. This experimental evidence highlights the essential role played by orthogonality constraints within OFT.

\subsection{General Framework}

Traditional finetuning generally relies on gradient descent using a modest learning rate to update all or select layers within a model. The rationale behind employing a small learning rate is to promote only minor modifications relative to the pretrained parameters, effectively resulting in finetuning that minimizes $\thickmuskip=2mu \medmuskip=2mu \| \bm{M} - \bm{M}^0\|$, where $\bm{M}$ and $\bm{M}^0$ represent the weights of the finetuned and pretrained models, respectively. While this implicit constraint appears reasonable, it may provide excessive freedom when adapting very large models. To enforce a more structured adaptation, LoRA imposes a low-rank restriction on the parameter update: $\thickmuskip=2mu \medmuskip=2mu \text{rank}(\bm{M} - \bm{M}^0)=r'$, with $r'$ chosen to be a small value. In contrast to LoRA, OFT constrains the pairwise similarity of neurons by enforcing $\thickmuskip=2mu \medmuskip=2mu \| \text{HE}(\bm{M})-\text{HE}(\bm{M}^0) \|=0$, where $\text{HE}(\cdot)$ is the hyperspherical energy associated with a given weight matrix. To make this more concrete, let us examine a fully connected layer parameterized as $\thickmuskip=2mu \medmuskip=2mu \bm{W}=\{\bm{w}_1,\cdots,\bm{w}_n\}\in\mathbb{R}^{d\times n}$, with $\bm{w}_i\in\mathbb{R}^{d}$ designating the $i$th neuron, and $\bm{W}^0$ the original, pretrained weights. The output for this layer, denoted $\thickmuskip=2mu \medmuskip=2mu \bm{z}\in\mathbb{R}^n$, is computed as $\thickmuskip=2mu \medmuskip=2mu \bm{z}=\bm{W}^\top\bm{x}$, where $\thickmuskip=2mu \medmuskip=2mu \bm{x}\in\mathbb{R}^d$ is an input. The OFT approach may thus be interpreted as an attempt to minimize the discrepancy in hyperspherical energy between the finetuned and original model:

\begin{equation}\label{eq:oft_principle}
\footnotesize
\min_{\bm{W}} \left\| \text{HE}(\bm{W})-\text{HE}(\bm{W}^0)\right\|~~\Leftrightarrow~~\min_{\bm{W}} \bigg{\|} \sum_{i\neq j} \|\hat{\bm{w}}_i-\hat{\bm{w}}_j\|^{-1}-\sum_{i\neq j} \|\hat{\bm{w}}^0_i-\hat{\bm{w}}^0_j\|^{-1}\bigg{\|}
\end{equation}

Here, the normalized $i$th neuron is given by $\thickmuskip=2mu \medmuskip=2mu \hat{\bm{w}}_i:=\bm{w}_i/\|\bm{w}_i\|$, and hyperspherical energy for the weight matrix $\bm{W}$ is defined as $\thickmuskip=2mu \medmuskip=2mu \text{HE}(\bm{W}):= \sum_{i\neq j} \|\hat{\bm{w}}_i-\hat{\bm{w}}_j\|^{-1}$. It is readily verified that the optimal value for Eq.~\eqref{eq:oft_principle} is zero. This is achievable whenever $\bm{W}$ results from applying an orthogonal transformation (rotation or reflection) to $\bm{W}^0$, i.e., $\thickmuskip=2mu \medmuskip=2mu \bm{W}=\bm{R}\bm{W}^0$ with $\thickmuskip=2mu \medmuskip=2mu \bm{R}\in\mathbb{R}^{d\times d}$ orthogonal (with determinant $\pm1$ signifying a rotation or a reflection, respectively). The OFT paradigm is predicated on this observation, advocating that finetuning may be accomplished simply by learning orthogonal

\begin{equation}\label{eq:oft_general}
    \footnotesize
    \bm{z}=\bm{W}^\top\bm{x}=(\bm{R}\cdot\bm{W}^0)^\top\bm{x},~~~\text{s.t.}~\bm{R}^\top\bm{R}=\bm{R}\bm{R}^\top=\bm{I}
\end{equation}

Here, $\bm{W}$ stands for the OFT-finetuned parameter matrix, with $\bm{I}$ being the identity matrix. An illustration of OFT is provided in Figure~\ref{fig:block}. To ensure that OFT fine-tuning commences from the pretrained $\bm{W}^0$, mirroring the zero initialization approach in LoRA, it is necessary to set the initial orthogonal matrix $\bm{R}$ as the identity matrix. This initialization makes certain that the fine-tuning process starts from the pretrained parameters. Ensuring that $\bm{R}$ remains orthogonal can be accomplished by employing differentiable orthogonalization methods, such as those presented in \cite{liu2021orthogonal,lezcano2019cheap}. Further discussion will address strategies for maintaining orthogonality efficiently during training.

\subsection{Efficient Orthogonal Parameterization}\label{sect:eff_ortho}

Conventional orthogonalization procedures, for instance, the Gram-Schmidt process, although differentiable, can introduce significant computational overhead~\cite{liu2021orthogonal}. To enhance computational efficiency, we utilize Cayley parameterization to construct orthogonal matrices. Specifically, we define the orthogonal matrix via $\thickmuskip=2mu \medmuskip=2mu \bm{R}=(\bm{I}+\bm{Q})(I-\bm{Q})^{-1}$, with $\bm{Q}$ representing a skew-symmetric matrix ($\thickmuskip=2mu \medmuskip=2mu \bm{Q}=-\bm{Q}^\top$). This parameterization provides efficiency benefits, although it restricts $\bm{R}$ to be in the special orthogonal group, i.e., orthogonal matrices with determinant $1$. In empirical evaluations, we observe that this restriction does not negatively affect performance. However, when dimensionality $d$ grows, directly parameterizing $\bm{R}$ in this manner may still require substantial parameters. To mitigate this, we instead model $\bm{R}$ as a block-diagonal matrix with $r$ distinct blocks, yielding the structure:

\begin{equation}
    \footnotesize
    \bm{R}=\text{diag}(\bm{R}_1,\bm{R}_2,\cdots,\bm{R}_r)=\begin{bmatrix}
    \bm{R}_1\in \text{O}(\frac{d}{r}) &  &\\
    &\ddots & \\
    & & \bm{R}_r\in \text{O}(\frac{d}{r})
    \end{bmatrix}\in \text{O}(d)
\end{equation}

Here, $\text{O}(d)$ refers to the orthogonal group in $d$ dimensions, and we define $\thickmuskip=2mu \medmuskip=2mu \bm{R}\in\mathbb{R}^{d\times d}$ as well as $\thickmuskip=2mu \medmuskip=2mu \bm{R}_i\in\mathbb{R}^{d/r\times d/r}$ for each $i$ as orthogonal matrices. In the special case when $\thickmuskip=2mu \medmuskip=2mu r=1$, the block-diagonal structure reduces to an unconstrained full orthogonal matrix. A generic orthogonal matrix of size $\thickmuskip=2mu \medmuskip=2mu d\times d$ possesses $\thickmuskip=2mu \medmuskip=2mu d(d-1)/2$ free parameters, yielding a computational complexity of $\mathcal{O}(d^2)$. For an orthogonal matrix partitioned into $r$ blocks, the total number of parameters decreases to $\thickmuskip=2mu \medmuskip=2mu d(d/r-1)/2$, thus requiring complexity on the order of $\thickmuskip=2mu \medmuskip=2mu \mathcal{O}(d^2/r)$. An additional parameter reduction can be achieved by reusing the block matrix across all blocks; in this scenario, $\thickmuskip=2mu \medmuskip=2mu\bm{R}_i=\bm{R}_j$ for any $i\neq j$, shrinking the parameter complexity down to $\mathcal{O}(d^2/r^2)$. Importantly, these parameter-efficient approaches all preserve the orthogonality of the resultant matrix $\bm{R}$, thereby maintaining the hyperspherical energy.

To compare the parameter efficiency of OFT and LoRA, consider that LoRA, with rank parameter $r'$, contains $\thickmuskip=2mu \medmuskip=2mu r'(d+n)$ trainable parameters. If both $r$ and $r'$ are defined as fractions of the dimension $d$, specifically $\thickmuskip=2mu \medmuskip=2mu r=r'=\alpha d$ with $0<\alpha\leq 1$, then LoRA's parameter count scales as $\mathcal{O}(d^2+dn)$, whereas OFT’s scales as $\mathcal{O}(d)$. To highlight this contrast, let us examine a weight matrix of dimensions $\thickmuskip=2mu \medmuskip=2mu 128\times 128$. With $r'=8$ for LoRA, the total number of trainable parameters is $2,048$. For OFT with $r=8$ and without shared blocks, the parameter count drops to $960$.

\subsection{Constrained Orthogonal Finetuning}

OFT’s expressiveness can be further curtailed by ensuring that the adapted parameters remain within a limited distance from the original pretrained parameters. In Constrained Orthogonal Finetuning (COFT), the forward propagation is described as follows:

\begin{equation}\label{eq:coft_general}
    \footnotesize
    \bm{z}=\bm{W}^\top\bm{x}=(\bm{R}\cdot\bm{W}^0)^\top\bm{x},~~~\text{s.t.}~\bm{R}^\top\bm{R}=\bm{R}\bm{R}^\top=\bm{I},~\norm{\bm{R}-\bm{I}}\leq \epsilon
\end{equation}

Here, the conditions require orthogonality of $\bm{R}$ and constrain its deviation from the identity matrix to be less than or equal to $\epsilon$. The Cayley parameterization, introduced earlier in Section~\ref{sect:eff_ortho}, ensures the orthogonality requirement. However, imposing the $\epsilon$-deviation constraint within this parameterization is nontrivial. To better understand the Cayley transform, we can use the Neumann series to approximate $\thickmuskip=2mu \medmuskip=2mu \bm{R}=(\bm{I}+\bm{Q})(I-\bm{Q})^{-1}$, yielding the first-order approximation $\thickmuskip=2mu \medmuskip=2mu \bm{R}\approx\bm{I}+2\bm{Q}+\mathcal{

series converges in the operator norm). This allows us to embed the constraint $\thickmuskip=2mu \medmuskip=2mu\|\bm{R}-\bm{I}\|\leq\epsilon$ directly within the Cayley transform. In turn, this leads to an equivalent requirement of $\thickmuskip=2mu \medmuskip=2mu \|\bm{Q}-\bm{0}\|\leq\epsilon'$, where $\bm{0}$ stands for a matrix filled with zeros and $\epsilon'$ is a new tolerance hyperparameter, generally distinct from $\epsilon$. This updated constraint on the matrix $\bm{Q}$ is straightforwardly enforced using projected gradient descent. To ensure that the orthogonal matrix $\bm{R}$ starts at the identity, we set $\bm{Q}$ to be initially the zero matrix. COFT can thus be seen as enforcing two explicit constraints: the first keeps the hyperspherical energy difference small, and the second tightly limits divergence from the pretrained model. While many existing finetuning approaches implement the second constraint implicitly, COFT formalizes it explicitly. Despite OFT's strong empirical results, we have found that this explicit deviation restriction in COFT provides further improvements in finetuning stability. Figure~\ref{fig:eps_coft} demonstrates how different values of $\epsilon$ modulate COFT's behavior: increasing $\epsilon$ makes the COFT model behave more like the OFT variant, while decreasing $\epsilon$ causes it to revert toward the properties of the original pretrained diffusion model for text-to-image synthesis.

\subsection{Re-scaled Orthogonal Finetuning}

As a straightforward augmentation to the standard OFT, we introduce a learnable scaling parameter for each neuron. This addition is motivated by the observation that adjusting the scale of neuron outputs leaves the hyperspherical energy unaffected, since normalization nullifies magnitude changes. The proposed forward computation is expressed as: $\thickmuskip=2mu \medmuskip=2mu\bm{z}=(\bm{R}\bm{W}^0\bm{D})^\top\bm{x}$\footnote{\textbf{Errata}: The NeurIPS camera-ready manuscript incorrectly stated the forward propagation of re-scaled OFT as $\thickmuskip=2mu \medmuskip=2mu\bm{z}=(\bm{D}\bm{R}\bm{W}^0)^\top\bm{x}$. Nevertheless, the published implementation remains correct and thus the Appendix~\ref{app:rescaled} results are valid.}, where $\thickmuskip=2mu \medmuskip=2mu \bm{D}=\text{diag}(s_1,\cdots,s_n)\in\mathbb{R}^{n\times n}$ is a parameterized diagonal matrix whose entries $s_1,\cdots,s_n$ are all positive and trainable. Unlike in standard OFT, where only the orthogonal matrix $\bm{R}$ is optimized (see Eq.~\eqref{eq:oft_general}), the re-scaled version learns both $\bm{D}$ and $\bm{R}$ simultaneously. This extension makes OFT more adaptable, incurring only a minimal increase in parameter count. Nevertheless, we focus on baseline OFT in our primary experiments to specifically isolate the effects of orthogonal transformations, though experiments suggest that incorporating re-scaling typically yields additional gains (refer to Appendix~\ref{app:rescaled}).

\section{Intriguing Insights and Discussions}

\textbf{OFT exhibits architecture-independence}. The OFT approach is, in theory, applicable to any neural network architecture. In practice, LoRA is commonly employed to adapt Transformer attention weights~\cite{hulora2022}. To enable direct comparison, our experimental setup restricts OFT to finetune attention matrices alone. Beyond traditional fully connected structures, OFT is naturally compatible with convolutional layers. The block-diagonal formation of $\bm{R}$ is particularly meaningful for convolutions, unlike LoRA’s parameterization. For instance, if the number of blocks matches the number of input channels, each block modulates a separate neuron channel, analogous to parameterizing depth-wise convolution filters~\cite{chollet2017xception}. If, however, every block is identical within $\bm{R}$, the transformation applies an orthogonal operator uniformly across channels. Further exploratory experiments on OFT-based convolution layer adaptation are reported in Appendix~\ref{app:convolution}.

\textbf{Connection to LoRA}. LoRA incorporates a low-rank matrix to ensure the pretrained weight matrix retains its original information and is not excessively altered. Conversely, OFT enforces an orthogonal (full-rank) transformation on the pretrained weights, maintaining the integrity of the pretrained information by preventing its disruption. The forward computation in OFT can be expressed as $\thickmuskip=2mu \medmuskip=2mu \bm{z}=(\bm{R}\bm{W}^0)^\top\bm{x}=(\bm{W}^0+(\bm{R}-\bm{I})\bm{W}^0)^\top\bm{x}$, where $\thickmuskip=2mu \medmuskip=2mu (\bm{R}-\bm{I})\bm{W}^0$ serves a similar function to the low-rank update employed in LoRA. Assuming that $\bm{W}^0$ has full rank, OFT also yields a low-rank weight update whenever $\thickmuskip=2mu \medmuskip=2mu \bm{R}-\bm{I}$ is itself low-rank. Both approaches possess hyperparameters that control trainable parameter count: LoRA uses a rank parameter $r'$, and OFT leverages a diagonal block size parameter $r$. Intriguingly, LoRA and OFT exemplify two fundamentally different forms of parameter efficiency: LoRA capitalizes on low-rank structure to economize on learnable parameters, whereas OFT benefits from structured sparsity, namely block-diagonal orthogonality.

\textbf{Why OFT converges faster}. Examination of Figure~\ref{fig:toy_ae} reveals that the most direct means of effecting semantic changes is to alter neuron orientations—precisely the kind of modification OFT is devised to make. Moreover, OFT can be interpreted as adapting neurons on a smooth hyperspherical manifold, an approach known to facilitate a more favorable optimization landscape. This theoretical advantage is corroborated empirically in \cite{liu2021orthogonal}.

\textbf{Why not minimize hyperspherical energy}. A fundamental distinction from \cite{liu2021orthogonal} lies in our choice not to minimize hyperspherical energy. In the context of classification, redundancy-free neuron arrangements are preferable; minimizing hyperspherical energy distributes neurons uniformly on the hypersphere. Such uniformity, however, is not an appropriate objective for finetuning, as it risks erasing valuable information acquired during pretraining.

\textbf{Balancing flexibility and regularity in finetuning.} Our findings reveal a fundamental compromise between flexibility and regularity during finetuning. Traditional finetuning approaches offer maximal flexibility, yet they tend to suffer from instability and are prone to model collapse. In contrast, OFT—remarkably straightforward in design—achieves a desirable middle ground by conserving pairwise neuron angles, thus maintaining a balance between these two competing factors. The use of block-diagonal parameterization can further be interpreted as enforcing even stronger regularization on the orthogonal matrix.

\textbf{No inference-time computational burden.} Distinct from methods such as ControlNet, the OFT framework does not add any computational overhead at inference for the finetuned model. During inference, it suffices to pre-multiply the learned orthogonal matrix $\bm{R}$ with the pretrained weight matrix $\bm{W}^0$, resulting in a single equivalent weight matrix $\thickmuskip=2mu \medmuskip=2mu \bm{W}=\bm{R}\bm{W}^0$. Therefore, inference throughput remains identical to that of the original pretrained model.

\section{Experiments and Results}

\textbf{Experimental setup.} Our experiments utilize Stable Diffusion v1.5~\cite{rombach2022high} as the underlying pretrained text-to-image generator. For a consistent evaluation, we randomly sample generated outputs from every competing method. The pipeline for subject-driven synthesis generally adheres to the procedures described in DreamBooth~\cite{ruiz2023dreambooth}, while the procedures for controllable image synthesis align with ControlNet~\cite{zhang2023adding} and T2I-Adapter~\cite{mou2023t2i}. In order to enable a direct comparison with LoRA, OFT or COFT is constrained to operate on the same network layers as LoRA. Additional details and comprehensive results are reported in Appendix~\ref{app:settings}.

\subsection{Subject-driven Generation}

\textbf{Configuration.} For subject-driven image generation, DreamBooth~\cite{ruiz2023dreambooth} and LoRA~\cite{hulora2022} serve as comparative baselines. All approaches are trained using the loss function employed in DreamBooth. For both DreamBooth and LoRA, we closely follow the settings and hyperparameters recommended in their respective publications. Extended experimental findings are included in Appendix~\ref{app:settings},\ref{app:coft_oft},\ref{app:more_qualitative},\ref{app:failure}.

\textbf{Finetuning stability and convergence}. Our initial experiments assess both the stability and convergence speed during finetuning for DreamBooth, LoRA, OFT, and COFT, with the results summarized in Figure~\ref{fig:teaser} and Figure~\ref{fig:db_convergence}. It is evident that COFT and OFT maintain stable finetuning behavior for the diffusion model across the training process. At the 400-iteration mark, both DreamBooth and OFT variants successfully exercise control, whereas LoRA is unable to preserve the subject’s identity. Progressing to 2000 iterations, DreamBooth suffers from image collapse, and LoRA incorrectly synthesizes yellow fur instead of a yellow shirt. Conversely, OFT and COFT both sustain robust and steady control over the generation outputs. These findings affirm the rapid yet reliable convergence offered by our OFT methodology for subject-driven generation tasks. Notably, the resilience with respect to the number of finetuning iterations proves valuable, as it eliminates much of the need for extensive tuning across different subjects. With both OFT and COFT, one can conveniently select a higher iteration count without intensive hyperparameter optimization. Particularly with COFT and an appropriately chosen $\epsilon$, setting both the learning rate and the number of iterations becomes straightforward.

\textbf{Quantitative comparison}. Consistent with the evaluation protocol from \cite{ruiz2023dreambooth}, we undertake a quantitative analysis to assess subject fidelity (via DINO~\cite{caron2021emerging}, CLIP-I~\cite{radford2021learning}), text prompt fidelity (CLIP-T~\cite{radford2021learning}), and sample diversity (LPIPS~\cite{zhang2018unreasonable}). CLIP-I evaluates the mean pairwise cosine similarity between CLIP embeddings of real and generated images; DINO applies a comparable approach using ViT S/16 DINO representations. For CLIP-T, we compute the average cosine similarity between the text prompt and generated image embeddings. Diversity is assessed by measuring the average LPIPS cosine similarity among images generated from identical prompts and subjects. As shown in Table~\ref{table:dreambooth_final}, both COFT and OFT significantly surpass DreamBooth and LoRA on the DINO and CLIP-I benchmarks, while demonstrating comparable or slightly superior results in terms of prompt fidelity and image diversity. Each approach is evaluated by repeatedly finetuning the same pretrained model with 30 different random seeds to account for stochasticity. Overall, our results demonstrate that the OFT framework not only converges more reliably and stably but also achieves consistently higher final performance.

\textbf{Qualitative comparison}. To provide clearer insights into the advantages of OFT, we present a selection of randomly chosen samples for subject-driven generation, as illustrated in Figure~\ref{fig:dreambooth_final}. For consistency and fairness, all outputs are produced using the same fine-tuned model architecture for each compared approach, with no individual text prompt optimization tailored to specific outputs. Each approach is represented by the model checkpoint that yields the highest validation CLIP score. As observed from Figure~\ref{fig:dreambooth_final}, OFT and COFT both demonstrate strong preservation of semantic subject features, whereas LoRA frequently fails to retain subject identity (for example, LoRA does not preserve subject identity in the bowl instance). Furthermore, OFT and COFT exhibit greater precision in responding to text prompts, whereas DreamBooth—despite its ability to maintain subject identity—often struggles to adhere to textual guidance (notably in the first row, bowl scenario). These qualitative assessments highlight that OFT not only offers improved control via text prompts but also consistently maintains subject information. Additionally, the iteration count within OFT is robust, with OFT sustaining performance even when the number of iterations is increased, a property not shared by DreamBooth or LoRA. Additional visual results are available in Appendix~\ref{app:more_qualitative}. Further, Appendix~\ref{app:human} describes a human study that provides further evidence for the advantages of OFT.

\subsection{Controllable Generation}

\textbf{Settings}. For baseline comparisons, we adopt ControlNet~\cite{zhang2023adding}, T2I-Adapter~\cite{mou2023t2i}, and LoRA~\cite{hulora2022}. We investigate three demanding tasks involving controllable image synthesis: Canny edge-to-image (C2I) on the COCO dataset~\cite{lin2014microsoft}, segmentation map-to-image (S2I) on the ADE20K dataset~\cite{zhou2017scene}, and landmark-to-face (L2F) on the CelebA-HQ dataset~\cite{karras2018progressive,xia2021tedigan}. For each approach, Stable Diffusion (SD) v1.5 is fine-tuned on the respective datasets for 20 epochs. Further details and results can be found in Appendix~\ref{app:more_qualitative},~\ref{app:more_control_task}, and~\ref{app:failure}.

\textbf{Convergence}. We assess the rate of convergence for ControlNet, T2I-Adapter, LoRA, and COFT within the context of the L2F task, incorporating both quantitative and qualitative analyses. For our metric, we utilize the mean $\ell_2$ distance calculated between the control face landmarks and the predicted landmarks. Figure~\ref{fig:face_convergence} illustrates the face landmark error as the models are fine-tuned over a varying number of epochs. The data demonstrates that COFT and OFT reach lower errors much more rapidly. For instance, LoRA requires 20 epochs to match the performance level that OFT achieves by the 8-th epoch. It is noteworthy that both OFT and COFT maintain a comparable number of trainable parameters to LoRA (substantially fewer compared to ControlNet), yet their convergence is markedly quicker than alternative techniques. Furthermore, Figure~\ref{fig:teaser} substantiates the efficient convergence behavior of OFT: the example on the right highlights that OFT needs only 5\% of the ADE20K dataset to converge effectively, making it more data-efficient than both ControlNet and LoRA. Qualitative assessments prioritize comparisons among OFT, COFT, and ControlNet, as ControlNet attains landmark errors closest to those of OFT. According to the results shown in Figure~\ref{fig:face_convergence_qual}, both OFT and COFT demonstrate consistent convergence with generated face poses incrementally aligning to the target landmarks. In contrast, ControlNet displays a sudden onset of controllability after the 8-th epoch, aligning with the abrupt convergence pattern previously reported in \cite{zhang2023adding}. The reliable and smooth convergence of OFT facilitates its practical deployment, as the training dynamics are more straightforward and predictable, simplifying the search for optimal hyperparameters.

\textbf{Quantitative comparison}. To assess the effectiveness of controllable generation, we propose a control consistency metric. This metric operates by extracting the control signal from the output image and measuring its similarity to the original input signal. Specifically, for the C2I task, we report IoU and F1 scores. For S2I, we utilize mean IoU along with mean and overall accuracy. In the case of L2F, the mean $\ell_2$ distance between the predicted and ground-truth control landmarks is calculated. Further methodological details on these consistency metrics can be found in Appendix~\ref{app:settings}. Across all evaluated approaches, we adopt their optimal hyperparameter settings. Table~\ref{table:control_gen} reveals that OFT and COFT consistently provide superior and more precise control compared to alternatives. We notice that methods relying on adapters (such as T2I-Adapter and ControlNet) display slower convergence and ultimately deliver inferior results. While LoRA surpasses ControlNet in the S2I setting, it falls short in both the C2I and L2F tasks. Overall, LoRA's maximum attainable performance appears limited, despite thorough hyperparameter optimization. By contrast, OFT continues to improve as the number of learnable parameters increases, suggesting its potential has yet to plateau. Our quantitative evaluation framework for controllable generation represents a key contribution, enabling precise assessment of finetuned model control performance, bounded by the accuracy of the adopted off-the-shelf segmentation or detection models.

\textbf{Qualitative comparison}. We further conduct a qualitative evaluation of OFT and COFT against leading approaches such as ControlNet, T2I-Adapter, and LoRA. As depicted by the randomly sampled outputs in Figure~\ref{fig:control_final}, OFT and COFT consistently deliver images that are both highly realistic and faithful in detail, while also maintaining precise control over the generation process. Focusing on the S2I task, it is evident that LoRA does not succeed in aligning the output with the provided segmentation map, whereas ControlNet, OFT, and COFT effectively manage this form of control. OFT and COFT, unlike ControlNet, generate images with finer details and more plausible geometry, and they do so with substantially fewer model parameters. In the C2I setting, OFT and COFT exhibit strong performance in synthesizing semantically coherent images from minimal Canny edge cues, whereas T2I-Adapter and LoRA lag behind in quality. Regarding the L2F task, our approach demonstrates superior accuracy in pose conditioning, achieving the best correspondence between generated and target face poses, even under difficult angles. Across all three control scenarios, OFT and COFT consistently outperform baseline methods, underscoring the robustness and versatility of our OFT framework for controlled image synthesis. For a more thorough qualitative assessment, additional results across the three primary control tasks are available in Appendix~\ref{app:control_exp}; we also extend the demonstration to further control settings—such as dense pose to body synthesis, sketch-to-image, and depth-to-image conversion—in Appendix~\ref{app:more_control_task}.

\section{Concluding Remarks and Open Problems}

Prompted by the recognition that the angular relationships between neurons play a vital role in shaping visual semantics, this work introduces a straightforward yet robust finetuning strategy—orthogonal finetuning—for managing text-to-image diffusion models. The approach is applied to two primary text-to-image tasks: subject-driven generation and controllable generation. OFT exhibits superior controllability and greater finetuning stability relative to previous approaches, all while requiring a reduced number of finetuning parameters. Notably, OFT does not incur any extra inference costs, thereby yielding a model well-suited for efficient deployment.

Nevertheless, OFT raises several noteworthy open questions. First, OFT enforces orthogonality through Cayley parametrization, which necessitates a matrix inversion and consequently can hinder the scalability of OFT. While the adoption of block diagonal parametrization offers a practical workaround, developing faster and differentiable methods for matrix inversion remains an open challenge. Second, OFT is distinctive in its potential for compositionality: orthogonal matrices obtained from separate OFT finetuning procedures can be combined via multiplication and the result remains orthogonal. Investigating whether this collection of orthogonal matrices successfully retains the information from all corresponding downstream tasks is an appealing avenue for future research. Lastly, the parameter efficiency achieved by OFT largely stems from the use of a block diagonal structure, a choice that inevitably introduces bias and constrains flexibility. Finding ways to enhance parameter efficiency while minimizing bias and maximizing adaptability stands out as a significant open research problem.